\documentclass{situationist-preprint}

\title{Against Illusionism About Free Will}
\msid{ACT-18-a}
\msfamily{ACT-18}
\msdiscipline{Action and Free Will}
\mstype{Principal manuscript}
\mskeywords{illusionism; Smilansky; free will; ultimacy; belief in free will}

\begin{document}
\maketitle

\begin{abstract}
Illusionism holds that we lack free will of the kind our moral practices presuppose, but that the belief that we have it is necessary, or at least beneficial, and should be sustained. We argue that illusionism is unmotivated. Its case has two parts: a metaphysical part, that the free will our practices presuppose does not exist, and a practical part, that the belief in it does valuable work. If free will is real, as ACT-01 argues, the practical part needs no illusion, since the belief that does the work is true. The illusionist's reply is that the free will that exists is compatibilist and shallow, while the free will our practices presuppose is libertarian and requires ultimacy. We argue that this reply fails: responsibility does not require ultimacy (ACT-08), and no way the world could be would supply the deeper freedom the illusionist describes.
\end{abstract}

\section{The position}

\citet{Smilansky2000} holds that libertarian free will is incoherent, that compatibilist free will exists but is shallow, and that the moral and personal goods of our practices of responsibility depend on our believing that we have something deeper. Since the truth would erode those goods if widely appreciated, he concludes that illusion plays, and should continue to play, a positive role. Hard incompatibilists, by contrast, hold that we can live well without the belief \citep{Pereboom2001}.

Illusionism needs two premises.

\begin{principle}\label{p:meta}
The free will our practices of responsibility presuppose does not exist.
\end{principle}

\begin{principle}\label{p:prac}
The belief that we have that free will does work that the practices need and that no true belief would do as well.
\end{principle}

Principle~\ref{p:prac} alone does not motivate illusionism, since a belief that is true and useful is not an illusion.

\section{The argument}

ACT-01-a argues that we have free will, understood as the capacity to act from a mechanism of practical reasoning that is one's own and that responds to reasons, and that this capacity is compatible both with determinism and with indeterminism. If that is the free will our practices presuppose, Principle~\ref{p:meta} is false, and illusionism is unmotivated.

The illusionist denies the antecedent. What our practices presuppose, on her view, is not reasons-responsiveness but ultimacy: that the agent be the ultimate source of what she does, so that nothing she did not choose explains it. Compatibilist freedom is real, but it leaves agents at the mercy of their unchosen natures, and the sense of this is what makes the compatibilist account seem shallow \citep{Smilansky2000}.

This reply fails at two points. First, our practices do not presuppose ultimacy. ACT-08-a argues that Strawson's Basic Argument shows ultimacy to be impossible but does not show it to be required: an agent who acts from an owned, reasons-responsive mechanism is responsible for what she does, and her responsibility for her nature accrues by degrees as her actions shape it. The reactive attitudes respond to the quality of will an action expresses \citep{Strawson1962}, and that is a feature of the agent's present mechanism, not of the origin of her initial nature. The belief that does work in the practices, that agents act from their own reasons and could respond to reasons to act otherwise, is the belief that ACT-01-a shows to be true.

Second, the illusionist's contrast between shallow and deep freedom presupposes that the deep kind would be available if the world were otherwise. But ACT-01-a holds that what threatens free will is neither determinism nor indeterminism. Undetermined self-forming choices, for which ACT-15-a argues agents can be responsible, still issue from efforts grounded in natures the agent did not choose, and so do not deliver ultimacy either. There is no way the world could be in which the deep freedom would exist, and practices that needed a belief in it would be incoherent rather than merely unfounded.

\begin{proposition}\label{p:main}
If free will is real in the sense of ACT-01, and our practices of responsibility presuppose no more than that sense, then Principle~\ref{p:meta} is false and illusionism is unmotivated.
\end{proposition}

\begin{proof}
The antecedents entail the negation of Principle~\ref{p:meta}, and a true belief is not an illusion.
\end{proof}

\section{An objection}

\begin{objection}[The empirical evidence]
Experiments suggest that weakening people's belief in free will makes them more likely to cheat \citep{VohsSchooler2008}. The belief is therefore valuable, whatever its truth.
\end{objection}

\begin{reply}
The finding, whatever its robustness, supports Principle~\ref{p:prac}, which we need not deny. It does nothing for Principle~\ref{p:meta}. If the belief is valuable and true, the conclusion to draw is that it should not be undermined by arguments that misdescribe what it is a belief in, not that it should be protected as an illusion.
\end{reply}

\section{Conclusion}

If free will is real, as ACT-01 argues, and responsibility does not require ultimacy, as ACT-08 argues, the belief on which the practices depend is true, and there is no illusion to sustain. Free will is not a necessary illusion, because it is not an illusion.

\begin{thebibliography}{9}\small
\bibitem[Pereboom(2001)]{Pereboom2001} Pereboom, D. (2001). \emph{Living Without Free Will}. Cambridge University Press.
\bibitem[Smilansky(2000)]{Smilansky2000} Smilansky, S. (2000). \emph{Free Will and Illusion}. Oxford University Press.
\bibitem[Strawson(1962)]{Strawson1962} Strawson, P.~F. (1962). Freedom and resentment. \emph{Proceedings of the British Academy}, 48, 1--25.
\bibitem[Vohs and Schooler(2008)]{VohsSchooler2008} Vohs, K.~D. and Schooler, J.~W. (2008). The value of believing in free will: Encouraging a belief in determinism increases cheating. \emph{Psychological Science}, 19(1), 49--54.
\end{thebibliography}

\end{document}
